\documentclass[tikz,border=6pt]{standalone}
\usepackage{peabra}
\begin{document}
\begin{tikzpicture}[x=1cm,y=1cm,
  ag/.style={doboz=27mm, font=\scriptsize, minimum height=13mm, inner sep=2pt},
  kul/.style={szurkedoboz=27mm, font=\scriptsize, minimum height=13mm, inner sep=2pt}]
  \node[sotetdoboz=58mm, minimum height=12mm] (alap) at (0,0) {IEC 61508\\{\scriptsize\mdseries alapszabvány az E/E/PE biztonsági rendszerekre}};
  \foreach \x/\t [count=\i] in {-8.4/{\textbf{Folyamatipar}\\IEC 61511\\{\tiny SIL 1--4}},-5.6/{\textbf{Gépipar}\\IEC 62061\\{\tiny SIL 1--3}},-2.8/{\textbf{Vasút}\\EN 50126, 50129,\\EN 50716 {\tiny SIL 1--4}},0/{\textbf{Atomenergia}\\IEC 61513,\\IEC 60880},2.8/{\textbf{Orvostechnika}\\IEC 60601,\\IEC 62304},5.6/{\textbf{Közúti jármű}\\ISO 26262\\{\tiny ASIL A--D}},8.4/{\textbf{Folyamat-, gép-}\\{\textbf{elemek}}\\IEC 61800-5-2 stb.}} {
    \node[ag] (s\i) at (\x,-2.3) {\t};
    \draw[nyilv] (alap.south) -- (s\i.north);
  }
  \node[font=\scriptsize\bfseries, text=peSzurke, anchor=west] at (-9.8,-3.75) {Az IEC 61508-tól független, saját fejlődésű szabványok:};
  \foreach \x/\t in {-8.4/{\textbf{Gépbiztonság}\\ISO 13849-1\\{\tiny PL a--e}},-5.6/{\textbf{Polgári repülés}\\ARP4754B, ARP4761A\\{\tiny DAL A--E}},-2.8/{\textbf{Repülési szoftver}\\DO-178C,\\DO-254},0/{\textbf{Orvosi kockázat}\\ISO 14971},2.8/{\textbf{Rendeltetésszerű}\\{\textbf{funkció}}\\ISO 21448 (SOTIF)},5.6/{\textbf{MI a járműben}\\ISO/PAS 8800,\\ISO/TS 5083},8.4/{\textbf{Kiberbiztonság}\\ISO/SAE 21434,\\IEC 62443}}
    \node[kul] at (\x,-4.95) {\t};
  \node[szurkedoboz=60mm, font=\scriptsize] at (0,1.55) {Közös alapok: ISO/IEC Guide 51, ISO 31000, IEC 31010, ISO 12100};
  \draw[nyilv] (0,1.15) -- (alap.north);
\end{tikzpicture}
\end{document}
